\section{Angolo di deviazione minima}
Quando un raggio luminoso attraversa un prisma vicino a un suo spigolo, subisce una doppia rifrazione nel passaggio dall'aria al prisma e successivamente dal prisma all'aria.
L'angolo formato tra il raggio incidente e quello uscente dal prisma viene indicato come angolo di deviazione totale. Tale angolo varia in base all'angolo di incidenza e assume una posizione limite, ovvero un valore minimo detto angolo di deviazione minima.
A partire dalla misura dell'angolo di deviazione minima si può ricavare l'indice di rifrazione del materiale tramite la relazione \\
\begin{equation}
n = \frac{\sin\!\left(\frac{\delta_{\mathrm{min}}+\alpha}{2}\right)}{\sin\!\left(\frac{\alpha}{2}\right)}
\label{angolo_deviazione_minima}
\end{equation}

dove $\alpha$ è l'angolo dello spigolo al vertice del prisma e $\delta_{\mathrm{min}}$ è l'angolo di deviazione minima.
Le misure dell'angolo di deviazione minima sono state effettuate sia per un prisma a base triangolare equilatera, sia per un prisma a base trapezoidale in corrispondenza dello spigolo $\alpha=45^\circ$.
Si riportano di seguito i dati raccolti insieme ai due indici di rifrazione calcolati tramite la relazione \eqref{angolo_deviazione_minima}.\\

\begin{table}[htbp]
\centering
\begin{tabular}{lcccc}
\toprule
\textbf{Prisma} & $\bm{\delta_{\text{min}}}$ \textbf{\small{(\unit{\degree})}} & $\bm{\alpha}$ \textbf{\small{(\unit{\degree})}} & $\mathbf{n}$ & $\bm{\sigma_{\text{n}}}$ \\
\midrule
A base triangolare  & 45 & 60 & 1.59 & 0.02 \\
A base trapezoidale & 27 & 45 & 1.54 & 0.02 \\
\bottomrule
\end{tabular}
\caption{Dati raccolti ed elaborati per l'angolo di deviazione minima.}
\label{tab:deviazione_minima}
\end{table}

Gli errori sugli indici di rifrazione indicati nella tabella sono dovuti alla propagazione dell'incertezza di $\ang{1}$  sulla misura degli angoli di incidenza e di deviazione minima.
L'indice di rifrazione del prisma a base trapezoidale, in plexiglass, risulta compatibile entro due deviazioni standard con entrambe le misure effettuate nella sezione precedente. Dalla relazione \eqref{eq:t-test} si calcola infatti $t_0=1.05$ con la \hyperref[misura-1]{prima misura}, $t_0=1.12$ con la \hyperref[misura-2]{seconda}. Per quanto riguarda invece il prisma a base triangolare, in vetro, non è possibile stabilire un confronto puntuale, poiché l'indice di rifrazione di questo materiale dipende dalla composizione chimica, non nota a priori. Tuttavia, la misura risulta pienamente compatibile con l'intervallo tipico di variabilità, indicato tra 1.50 e 1.90 \cite{wikipedia_vetro}.
